
\documentclass[12pt,a4paper]{article}
\usepackage[utf8]{inputenc}
\usepackage[portuguese]{babel}
\usepackage{booktabs}
\usepackage{geometry}
\usepackage{amsmath}
\usepackage{float}

\geometry{margin=2.5cm}

\title{Análise de Desempenho: Regressores Supervisionados}
\author{Estudante UNIFOR}
\date{\today}

\begin{document}

\maketitle

\section{Métricas de Avaliação}
Para a avaliação dos algoritmos de regressão, foram utilizadas métricas que medem a aderência do modelo aos dados contínuos. Conforme estabelecido no projeto, o foco reside na capacidade de explicação da variância:

\begin{itemize}
    \item \textbf{R2-Score (Coeficiente de Determinação):} Mede a proporção da variância dos dados que é explicada pelo modelo. Varia de $-\infty$ a 1.
    \item \textbf{R2 Ajustado:} Uma variação do R2 que penaliza a inclusão de atributos que não contribuem significativamente para o poder preditivo, sendo essencial para modelos com múltiplas variáveis.
\end{itemize}

Importante notar que, embora o roteiro cite métricas como Acurácia e F1-Score para regressores, estas não foram aplicadas. Em problemas de regressão (valores contínuos), a probabilidade de prever o valor exato é nula, tornando métricas de classificação tecnicamente inadequadas para esta tarefa específica.

\section{Resultados}
A tabela abaixo apresenta o desempenho médio, desvio padrão e tempos de execução para os três modelos implementados, utilizando validação cruzada (K-Fold) e limitação de 10.000 amostras para viabilidade computacional.

\begin{table}[H]
\centering
\caption{Análise comparativa do desempenho dos regressores.}
\label{tab:resultados}
\begin{tabular}{lcccc}
\toprule
\textbf{Algoritmo} & \textbf{R2-Score} & \textbf{R2 Ajustado} & \textbf{T. Treino (s)} & \textbf{T. Teste (s)} \\
\midrule
kNN (Euclidiana) & $0.1724 \pm 0.012$ & $0.1710 \pm 0.012$ & $0.0000 \pm 0.000$ & $679.96 \pm 2.03$ \\
kNN (Manhattan)  & $0.2009 \pm 0.012$ & $0.1995 \pm 0.012$ & $0.0000 \pm 0.000$ & $642.08 \pm 3.38$ \\
Reg. Linear Múltipla & $0.3014 \pm 0.024$ & $0.2972 \pm 0.024$ & $0.0004$ & $0.0001$ \\
\bottomrule
\end{tabular}
\end{table}

\section{Conclusões}
Com base nos experimentos realizados, a \textbf{Regressão Linear Múltipla} apresentou o melhor equilíbrio entre desempenho e eficiência computacional.

\begin{enumerate}
    \item \textbf{Desempenho:} A Regressão Linear obteve o maior R2-Score (0.3014), superando ambas as variações do kNN. Isso indica que a relação entre as variáveis do dataset \textit{interest\_rate} possui componentes lineares que foram melhor capturados por este modelo.
    \item \textbf{Eficiência:} O kNN (preguiçoso por natureza) demonstrou um custo computacional extremamente elevado na fase de teste (mais de 10 minutos por fold), enquanto a Regressão Linear, baseada na Equação Normal, executou predições em frações de milissegundo.
    \item \textbf{Limitação:} O baixo valor geral de R2 em todos os modelos sugere que o problema é complexo e que os atributos numéricos disponíveis podem não ser suficientes para explicar a totalidade da variação da taxa de juros sem engenharia de atributos adicional.
\end{enumerate}

\end{document}
